% !TEX root = chapter-standalone.tex

\section{Stacking categories}
\label{sec:parallel-stacking}

\todotext{This is strange here, move to intro?}
So far we have seen how we can compose morphisms ``horizontally'':
%
\begin{equation}
    \prfperiod{
        \mora \colon \Obja \mto \Objb
    }{
        \morb\colon \Objb \mto \Objc
    }{
        \morab \colon \Obja \mto \Objc
    }
\end{equation}
%
There are other notions of composition that allow us to compose morphisms by ``stacking them vertically''.
Given two morphisms
%
\begin{align}
    \mora & \colon \Obja \mto \Objb, \\
    \morb & \colon \Objc \mto \Objd,
\end{align}
%
we will obtain by parallel stacking a morphism
%
\begin{align}
    (\mora\mtimescatmor \morb)
     & \colon (\Obja \mtimescatob \Objc) \mto (\Objb \mtimescatob \Objd),
\end{align}
%
where ``$\mtimescatmor$'' and ``$\mtimescatob$'' are operations to be defined.
Note that while in the case of sequential composition there was a compatibility condition to be defined, as the target of the first morphism must be the source of the second morphism, here instead we can stack arbitrary morphisms.

We also expect to be able to stack any number of morphisms.
Having a collection of morphisms
%
\begin{equation}
    \moran{i} \colon \Objan{i} \mto \Objbn{i}, \qquad 1 \leq i \leq n,
\end{equation}
%
we expect to be able to obtain the composed morphism
%
\begin{equation}
    (\mtimescatmor_{i=1}^n \moran{i} )
    \colon( \mtimescatob_{i=1}^n \Objan{i} )\mto (\mtimescatob_{i=1}^n \Objbn{i}).
\end{equation}


\begin{marginfigure}
    \centering
    \includesag{parallelism-stacked-morphisms}
    %\includegraphics[scale=0.12]{parallelism_stacked-morphisms}
    \caption{Stacked morphisms}
    \label{fig:stacked-morphisms}
\end{marginfigure}

\begin{marginfigure}
    \centering
    \includesag{parallelism_stacked-morphisms-diagram}
    %\includegraphics[scale=0.17]{parallelism_stacked-morphisms-diagram}
    \caption{Stacking string diagrams}
    \label{fig:stacked-morphisms-diagram}
\end{marginfigure}

\begin{ctdefinition}[Stacking category]
    \label{def:simple-stacking-cat}
    A \maindef{stacking category} is a \SY{category}~\CatC with the following additional constituents and properties.

    \constit
    \begin{itemize}
        \item A stacking operation~$\mtimescatob \colon \ObC \cartprod \ObC \sto \ObC$.
        \item A stacking operation~$\mtimescatmor \colon \Mor_\CatC \cartprod \Mor_\CatC \sto \Mor_\CatC$.
    \end{itemize}

    \condit
    \begin{itemize}
        \item The two operations~$\mtimescatob$ and~$\mtimescatmor$ are compatible in the sense that
              %
              \begin{equation}
                  \prfperiod{
                      \moran{1} \colon \Objan{1} \mto \Objbn{1}
                      \qquad
                      \moran{2} \colon \Objan{2} \mto \Objbn{2}
                  }{
                      \moran{1} \mtimescatmor \moran{2} \colon  \Objan{1} \mtimescatob  \Objan{2} \mto \Objbn{1} \mtimescatob \Objbn{2}
                  }
              \end{equation}
    \end{itemize}
\end{ctdefinition}

In \cref{fig:stacked-morphisms} we have depicted a string diagram of two stacked morphisms.
Alternatively, in \cref{fig:stacked-morphisms-diagram} we depict the stacking of the string diagrams for $\moran{1}$ and $\moran{2}$, respectively, by stacking their diagrams vertically and drawing a box around them, merging their respective input and output terminals.
The outer box denotes $\moran{1} \mtimescatmor \moran{2}$; we think of \cref{fig:stacked-morphisms} as a ``black-boxed'' version of \cref{fig:stacked-morphisms-diagram}.

\begin{figure*}[tbh]
    \centering
    \begin{subfigure}{0.49\linewidth}
        \includesag{parallelism_associative-stacking-1}
        %\includegraphics[scale=0.15]{parallelism_associative-stacking-1}
        \caption{One way of stacking three morphisms}
        \label{fig:associative-stacking-1}
    \end{subfigure}
    %
    \begin{subfigure}{0.49\linewidth}
        \includesag{parallelism_associative-stacking-2}
        %\includegraphics[scale=0.15]{parallelism_associative-stacking-2}
        \caption{Another way of stacking three morphisms}
        \label{fig:associative-stacking-2}
    \end{subfigure}
    \caption{Stacking three morphisms.}
\end{figure*}

\todotext{Put examples here of stacking categories that are in the slides}


\toexclude{
    \begin{example}
        After the X101 spontaneously combusted in low Earth orbit, the astronauts at Jeb's Spaceship Parts go on strike.
        They demand that the engineers take into account safety and living conditions on the future X102.
        As long as the propulsion and life support systems of the X102 do not interact, we can simply tensor the two \SY{design problems} representing these systems into one, big co-design problem (\cref{fig:examplemonoidal}).
        \begin{figure}[h!]
            \centering
            \includesag{50_engine_tensor_1_2}
            \caption{Example of tensor of \SY{design problems}. }
            \label{fig:examplemonoidal}
        \end{figure}
    \end{example}
}


